%%% FIELD proof-logged-enclosure
Put \(t_0=\lceil n\varepsilon_n\rceil\).  On \(G_n\), the first component in the definition of \(G_n\) gives
\[
 \left|\bar q(t/n)-t e_{nt}\right|\leq\delta_{q,n},
 \qquad t_0\leq t\leq n.
\]
The three defining properties of \(\mathcal Q\) follow respectively from
\(\mathrm{ass:intensity\mbox{-}anchor}\),
\(\mathrm{ass:intensity\mbox{-}bounds}\), and
\(\mathrm{ass:intensity\mbox{-}lipschitz}\).  Hence the displayed inequalities are exactly the remaining membership constraints in \(\mathcal Q_n(e)\), proving
\(\bar q\in\mathcal Q_n(e)\) on \(G_n\).

Fix \(q'\in\mathcal Q_n(e)\).  At every constrained grid point,
\[
 |q'(t/n)-\bar q(t/n)|
 \leq |q'(t/n)-t e_{nt}|+|t e_{nt}-\bar q(t/n)|
 \leq 2\delta_{q,n}.
\]
If \(u\geq t_0/n\), choose a constrained grid point \(t/n\) with
\(|u-t/n|\leq n^{-1}\).  The two Lipschitz inequalities yield
\[
 |q'(u)-\bar q(u)|
 \leq 2\delta_{q,n}+2L_q/n.
\]
If \(u<t_0/n\), both profiles equal \(c\) at zero, and therefore
\[
 |q'(u)-\bar q(u)|
 \leq |q'(u)-c|+|\bar q(u)-c|
 \leq 2L_q u
 \leq 2L_q(\varepsilon_n+n^{-1}).
\]
Both bounds are at most
\(2\delta_{q,n}+2L_q(\varepsilon_n+n^{-1})=d_{Q,n}\), which proves the asserted uniform enclosure.  No conditioning on the completed log is used.

%%% FIELD proof-bandwidth-vanishing
For every real \(x\), \(x\leq 2^{-n}\lceil2^n x\rceil<x+2^{-n}\).  Applying this deterministic inequality to \(x_n\) gives
\[
 x_n\leq h_n\leq x_n+2^{-n}.
\]
It remains to verify \(x_n\to0\) from the primitive sequences.  Write
\(\varepsilon_n^\star=\max\{\varepsilon_n,\delta_{q,n},n^{-1/2}\}\).
Under the corrected cutoff \(T_n\leq n+1\), one has \(0<\ell_n\leq1\);
since \(\varepsilon_n^\star\to0\), the cap is inactive eventually and
\[
 \ell_n=\frac{T_n-1}{n}\leq
 \varepsilon_n^\star+n^{-1},\qquad
 \ell_n\geq \varepsilon_n^\star-n^{-1}.
\]
In particular, eventually \(\ell_n\geq (2\sqrt n)^{-1}\), so
\(\log(1/\ell_n)=O(\log n)\) and
\((n\ell_n)^{-1}=O(n^{-1/2})\).

Let \(t_0=\lceil n\varepsilon_n\rceil\).  Since
\(\sum_{t=t_0}^{T_n-1}t\leq T_n^2/2\),
\[
 v_{E,n}\leq v_{0,n}+\frac{2B^2}{c}(T_n/n)^2\longrightarrow0.
\]
Also \(v_{C,n}=O(n^{-1})\) and
\(v_{Z,n}\leq B^2(\varepsilon_n^\star+n^{-1})^2/c\to0\).
The definitions and \(\delta_{r,n},\delta_{q,n}\to0\) give
\[
 D_{\mathrm{mark},n}\to0,\qquad d_{Q,n}\to0.
\]
Moreover,
\[
 0\leq D_{\mathrm{grid},n}
 \leq n^{-1}\{L_jK\,O(\log n)+L_r\}\longrightarrow0.
\]
Every summand in the definition of \(x_n\) therefore tends to zero, uniformly because all constants and modulus sequences are common to the declared class.  Thus \(x_n\to0\), and the dyadic bracket implies \(h_n\to0\).  Finally,
\[
 0\leq5(h_n-x_n)\leq5\cdot2^{-n},
\]
which is the claimed deterministic radius charge.

%%% FIELD proof-finite-transfer
Suppress the first index \(n\), put \(u_t=t/n\), \(p_t=\bar q(u_t)/t\), and
\[
 j(u)=\frac{u r(u)}{\bar q(u)}.
\]
The restrictions \(t\geq T_n>K\geq\bar q(u_t)\) imply \(0<p_t<1\).
We construct all comparisons sequentially on an extension of the original randomization space.

\emph{Predictable Bernoulli comparison.}
For two Bernoulli probabilities \(e,p\), the common-uniform construction has mismatch probability \(|e-p|\).  Applying it at time \(t\), conditional on \(\mathcal F_{n,t-1}\), and using fresh independent uniforms gives variables \(B_t\) such that the \(B_t\)'s are independent with
\(\Pr(B_t=1)=p_t\), while
\[
 \Pr(A_{nt}\ne B_t\mid\mathcal F_{n,t-1})=|e_{nt}-p_t|.
\]
Let \(\sigma_q\) be the first \(t\geq T_n\) for which
\(|t e_{nt}-\bar q(u_t)|>\delta_{q,n}\), with value \(n+1\) if there is no such time.  This is a predictable stopping rule because \(e_{nt}\) is \(\mathcal F_{n,t-1}\)-measurable and \(\bar q\) is deterministic.  Before \(\sigma_q\),
\[
 |e_{nt}-p_t|\leq\delta_{q,n}/t.
\]
Consequently, by the union bound and iterated expectation,
\[
 \Pr\{\exists t\geq T_n:t<\sigma_q,\ A_{nt}\ne B_t\}
 \leq\delta_{q,n}\sum_{t=T_n}^n t^{-1}
 =\delta_{q,n}H_{1,n}.
\]
On \(G_n\), \(\sigma_q=n+1\).

\emph{Replacement of predictable marks.}
On \(G_n\), for \(t\geq T_n\) the inequalities
\(\bar q(u_t)\geq c/2\) and
\(\delta_{q,n}\leq c/4\) give \(t e_{nt}\geq c/4\).
Using \(|r_{nt}(1)-r(u_t)|\leq\delta_{r,n}\) and
\(|r(u_t)|\leq B\),
\[
 \left|\frac{r_{nt}(1)}{n e_{nt}}-j(u_t)\right|
 \leq u_t\left\{\frac{4\delta_{r,n}}{c}
              +\frac{8B\delta_{q,n}}{c^2}\right\}.
\]
On \(G_n\) and on the no-mismatch event, the difference between the late rare-arm martingale and
\(\sum_{t=T_n}^n(B_t-p_t)j(u_t)\) is bounded by
\[
 \left\{\frac{4\delta_{r,n}}{c}
              +\frac{8B\delta_{q,n}}{c^2}\right\}
       \sum_{t=T_n}^nB_tu_t+\delta_{r,n},
\]
because \(p_tj(u_t)=r(u_t)/n\).  Since
\[
 \mathbb E\sum_{t=T_n}^n B_tu_t
 =\frac1n\sum_{t=T_n}^n\bar q(u_t)\leq K,
\]
the expectation of this nonnegative error envelope is at most
\[
 (1+4K/c)\delta_{r,n}+(8KB/c^2)\delta_{q,n}
 =D_{\mathrm{mark},n}.
\]

\emph{Bernoulli-to-Poisson comparison.}
For \(0\leq p\leq1\), direct comparison of the masses of
\(\operatorname{Bernoulli}(p)\) and \(\operatorname{Poisson}(p)\) gives total variation
\[
 p(1-e^{-p})\leq p^2.
\]
Indeed, the positive discrepancies of the Poisson law occur at zero and at counts at least two, and their sum is \(p(1-e^{-p})\).  Coupling independently coordinate by coordinate and taking a union bound therefore replaces the \(B_t\)'s by independent \(P_t\sim\operatorname{Poisson}(p_t)\) with failure probability at most
\[
 \sum_{t=T_n}^n p_t^2\leq K^2H_{2,n}\leq2K^2H_{2,n}.
\]

\emph{Transfer from the grid to the Poisson random measure.}
For \(I_t=((t-1)/n,t/n]\), put
\(\lambda_t=\int_{I_t}\bar q(v)\,dv/v\).  For \(v\in I_t\),
\[
 \left|\frac{\bar q(v)}v-\frac{\bar q(u_t)}{u_t}\right|
 \leq \frac1n\left(\frac{L_q}{v}+\frac K{v^2}\right).
\]
Integration and summation over \((\ell_n,1]\) yield
\[
 \sum_{t=T_n}^n|\lambda_t-p_t|
 \leq n^{-1}\{L_q\log(1/\ell_n)+K(1/\ell_n-1)\}
 =I_{\mathrm{grid},n}.
\]
Two Poisson variables of means \(\lambda,\mu\) can be constructed from a common
\(\operatorname{Poisson}(\min\{\lambda,\mu\})\) count and independent excess counts; the probability of unequal counts is
\(1-e^{-|\lambda-\mu|}\leq|\lambda-\mu|\).
Thus the count vectors can be coupled with failure probability at most
\(I_{\mathrm{grid},n}\).  Conditional on the matched cell count, place the points with density proportional to \(\bar q(v)/v\) in that cell.

For \(u,v\in[0,1]\), the bounds on \(\bar q,r\) and their Lipschitz moduli give
\[
 |j(u)-j(v)|
 \leq\left\{\frac{2(B+L_r)}c+\frac{4BL_q}{c^2}\right\}|u-v|
 =L_j|u-v|.
\]
Therefore the expected absolute mark-location error is at most
\[
 \sum_{t=T_n}^n\int_{I_t}L_j|u_t-v|\frac{\bar q(v)}v\,dv
 \leq \frac{L_jK}{n}\log(1/\ell_n).
\]
The right-endpoint Riemann bound for the compensating drift is
\[
 \left|\frac1n\sum_{t=T_n}^nr(u_t)-\int_{\ell_n}^1r(v)\,dv\right|
 \leq L_r/n.
\]
Hence the grid mark and drift error has expectation at most
\(D_{\mathrm{grid},n}\).  Markov's inequality, proved by
\(\mathbb E X\geq h\Pr(X>h)\) for \(X\geq0\), allocates failure probability
\((D_{\mathrm{mark},n}+D_{\mathrm{grid},n})/h_n\) to the combined mark and grid error exceeding \(h_n\).

\emph{Early rare-arm, control-arm, and omitted Poisson terms.}
Conditional orthogonality of martingale differences gives
\[
 \mathbb E S_{n,\mathrm{control}}^2
 =\mathbb E\sum_{t=1}^n
 \frac{e_{nt}r_{nt}(0)^2}{n^2(1-e_{nt})}
 \leq\frac{B^2\rho}{n(1-\rho)}
 =v_{C,n}.
\]
Stop the early rare-arm martingale immediately before an increment whose predictable variance would make its accumulated bracket exceed the declared budget.  The stopped martingale agrees with the original one on \(G_n\).  The part before
\(t_0=\lceil n\varepsilon_n\rceil\) has bracket at most \(v_{0,n}\) there.  For \(t_0\leq t<T_n\), stop also before the propensity-profile constraint fails.  On the retained path, \(e_{nt}\geq c/(4t)\), whence
\[
 e_{nt}(1-e_{nt})
 \left\{\frac{r_{nt}(1)}{n e_{nt}}\right\}^2
 \leq\frac{4B^2t}{cn^2}.
\]
Thus the stopped early rare-arm martingale has second moment at most \(v_{E,n}\).

For a Poisson random measure, the expectation and variance identities used here follow first for a simple integrand from independence and
\(\operatorname{Var}\{\operatorname{Poisson}(\lambda)\}=\lambda\), and then for the present integrands by \(L^1\) and \(L^2\) approximation on \((\epsilon,1]\) followed by \(\epsilon\downarrow0\).  In particular, the centered omitted integral on \((0,\ell_n]\) has variance
\[
 \int_0^{\ell_n}\frac{u^2r(u)^2}{\bar q(u)^2}
                 \frac{\bar q(u)}u\,du
 =\int_0^{\ell_n}\frac{u r(u)^2}{\bar q(u)}\,du
 \leq \frac{B^2\ell_n^2}{c}=v_{Z,n}.
\]
For any centered square-integrable \(X\),
\(\Pr(|X|>h_n)\leq\mathbb EX^2/h_n^2\), by applying the preceding Markov argument to \(X^2\).  A union bound therefore charges
\((v_{E,n}+v_{C,n}+v_{Z,n})/h_n^2\) for these three terms.  On the complement of all charged events, the mark-grid error and the three centered remainder terms are each at most \(h_n\); hence
\(|S_n-Z_{r,\bar q}|\leq4h_n\).  Adding \(\Pr(G_n^c)\leq\kappa_n\) and all coupling failures proves exactly
\[
 \Pr\!\left(G_n^c\cup\{|S_n-Z_{r,\bar q}|>4h_n\}\right)
 \leq D_{\mathrm{tr},n}.
\]

Finally, every displayed term is an explicit finite sum or elementary expression in the supplied bounds.  To prove uniform convergence, note that eventually
\[
 H_{1,n}\leq1+\log(1/\ell_n),\qquad
 H_{2,n}\leq(n\ell_n)^{-1}=O(n^{-1/2}),
\]
and
\[
 \delta_{q,n}\log(1/\ell_n)\to0,\qquad
 I_{\mathrm{grid},n}\to0.
\]
The first limit follows because \(\ell_n\) is eventually of order at least
\(\max\{\varepsilon_n,\delta_{q,n},n^{-1/2}\}\), using
\(\delta\log(1/\delta)\to0\) at zero and the assumed
\(\delta_{q,n}\log(1/\varepsilon_n)\to0\).
The estimates in the bandwidth proof give
\(D_{\mathrm{mark},n}+D_{\mathrm{grid},n}\to0\) and
\(V_n:=v_{E,n}+v_{C,n}+v_{Z,n}\to0\).  Since
\[
 h_n\geq\sqrt{D_{\mathrm{mark},n}+D_{\mathrm{grid},n}},
 \qquad h_n\geq V_n^{1/4},
\]
the two ratio terms are bounded by
\(\sqrt{D_{\mathrm{mark},n}+D_{\mathrm{grid},n}}\) and
\(\sqrt{V_n}\), respectively.  Together with \(\kappa_n\to0\), this proves
\(D_{\mathrm{tr},n}\to0\) uniformly over the declared class.

%%% FIELD proof-profile-continuity
Let \(d_r=\|r-r'\|_\infty\) and \(d_q=\|q-q'\|_\infty\).
Use a Poisson random measure on
\((0,1]\times[0,K]\) with intensity \(du\,dv/u\).  A point \((u,v)\) is retained for profile \(q\) when \(v\leq q(u)\), and for \(q'\) when \(v\leq q'(u)\).  This is a common construction of the two intensity measures.

At a point retained by both processes,
\[
 \left|\frac{ur(u)}{q(u)}-\frac{ur'(u)}{q'(u)}\right|
 \leq u\left(\frac{2d_r}{c}+\frac{4B d_q}{c^2}\right).
\]
Integrating this bound against the dominating intensity \(K\,du/u\) gives an expected common-point cost at most
\[
 \frac{2K}{c}d_r+\frac{4KB}{c^2}d_q.
\]
At a point retained by only one process, the total expected absolute mark cost from both one-sided excess regions is bounded by
\[
 \int_0^1\frac{2B}{c}|q(u)-q'(u)|\,du
 \leq\frac{2B}{c}d_q.
\]
These integrals are finite at zero.  Construct first on \((\epsilon,1]\) and let
\(\epsilon\downarrow0\); the expected omitted absolute raw-mark mass is at most
\(\int_0^\epsilon B\,du=B\epsilon\), so the coupled raw sums converge in \(L^1\).
The compensating drifts are \(\int_0^1r(u)\,du\) and
\(\int_0^1r'(u)\,du\), whose difference is at most \(d_r\).
Therefore this coupling satisfies
\[
 \mathbb E|Z_{r,q}-Z_{r',q'}|
 \leq\left(\frac{2K}{c}+1\right)d_r
 +\left(\frac{4KB}{c^2}+\frac{2B}{c}\right)d_q.
\]
Taking the infimum over couplings proves the asserted \(W_1\) bound.
Finally \(H_{b,h_n}\) is \(h_n^{-1}\)-Lipschitz, so for the displayed coupling
\[
 \left|\mathbb EH_{b,h_n}(Z_{r,q})
       -\mathbb EH_{b,h_n}(Z_{r',q'})\right|
 \leq h_n^{-1}\mathbb E|Z_{r,q}-Z_{r',q'}|,
\]
which proves the ramp-risk conclusion.

%%% FIELD proof-certified-coverage-oracle
This proof uses the corrected dyadic threshold search in
\(\mathrm{def:certificate\mbox{-}algorithm}\).

\emph{Finite termination.}
The pointwise bounds and adjacent Lipschitz constraints on a time mesh form a finite system of rational interval inequalities.  Subdivision of the value boxes and piecewise-linear interpolation therefore gives finite outer sup-norm nets of
\(\mathcal Q_n(e)\) and \(\mathcal R\) at any prescribed rational accuracy.
For a representative pair, truncate the Poisson integral on \((0,\gamma]\).
The centered omitted term has second moment at most
\[
 \int_0^\gamma\frac{u r(u)^2}{q(u)}\,du
 \leq B^2\gamma^2/c,
\]
and hence first absolute moment at most \(B\gamma/\sqrt c\).
On \([\gamma,1]\), the total Poisson intensity is at most
\(K\log(1/\gamma)<\infty\).  Quantize time and signed marks, truncate the total count, and bound the remainder by the elementary Poisson-series inequality
\[
 \Pr\{\operatorname{Poisson}(\Lambda)\geq m+1\}
 \leq\frac{\Lambda^{m+1}}
 {(m+1)!\{1-\Lambda/(m+2)\}},
 \qquad m+2>\Lambda.
\]
This follows by bounding every successive ratio in the tail series by
\(\Lambda/(m+2)\).  The continuity lemma charges the finite profile-net error.
Thus, by successively decreasing \(\gamma\) and the quantization meshes and increasing \(m\), all outward errors can be made no larger than the positive rational budget \(\zeta_n\).  Every remaining convolution is finite and uses directed rational arithmetic; logarithms are enclosed by their convergent rational series.

If a guard in the definition fires, the algorithm stops immediately.  Otherwise
\(\Delta_n<\alpha\).  For every \(q'\in\mathcal Q\) and \(r'\in\mathcal R\),
\[
 \operatorname{Var}(Z_{r',q'})
 =\int_0^1\frac{u r'(u)^2}{q'(u)}\,du
 \leq B^2/c.
\]
Consequently
\(\mathbb EH_{b,h_n}(Z_{r',q'})
\leq\Pr(|Z_{r',q'}|>b)\leq B^2/(cb^2)\).
For a sufficiently large dyadic grid point this strict upper bound is below
\(\alpha-\Delta_n\), so a finite outward calculation certifies some grid point.
The corrected increasing grid search therefore terminates and returns a well-defined \(a_n\).

\emph{Selection-safe finite coverage.}
For deterministic \(q\), define
\[
 F_q(b)=\sup_{r'\in\mathcal R}
          \mathbb EH_{b,h_n}(Z_{r',q}),\qquad
 A_q=\inf\{b\geq0:F_q(b)\leq\alpha-\Delta_n\}.
\]
The ramp is nonincreasing in \(b\) and \(h_n^{-1}\)-Lipschitz in \(b\);
therefore \(F_q\) is continuous and nonincreasing.  The preceding variance bound makes its feasible set nonempty, so the infimum is attained.
On \(G_n\), the logged-enclosure lemma gives
\(\bar q\in\mathcal Q_n(e)\).  For each realized completed log on this event, the certificate at \(a_n\) implies
\[
 F_{\bar q}(a_n)\leq\alpha-\Delta_n,
\]
and hence \(a_n\geq A_{\bar q}\).  Notice that \(A_{\bar q}\) is deterministic for the fixed experiment; no conditional distribution given the completed log has been invoked.

On the coupling supplied by the finite-transfer theorem,
\[
 \{G_n,\ |S_n|>a_n+5h_n,\ |S_n-Z_{r,\bar q}|\leq4h_n\}
 \subseteq\{|Z_{r,\bar q}|>A_{\bar q}+h_n\}.
\]
At the right-hand event the ramp \(H_{A_{\bar q},h_n}\) equals one.  Since the actual deterministic residual profile belongs to \(\mathcal R\),
\[
 \Pr(|Z_{r,\bar q}|>A_{\bar q}+h_n)
 \leq F_{\bar q}(A_{\bar q})
 \leq\alpha-\Delta_n.
\]
The transfer theorem now gives
\[
 \Pr(|\widehat\tau_n-\tau_n|>a_n+5h_n)
 \leq D_{\mathrm{tr},n}+\alpha-\Delta_n
 =\alpha-\zeta_n\leq\alpha.
\]
If a guard fires, \(C_n=[-2M,2M]\); and
\(|\tau_n|\leq2M\) by the bounded-outcome assumption.  Otherwise the last display proves coverage of the untruncated interval.  Intersecting it with
\([-2M,2M]\) cannot remove \(\tau_n\), proving
\(\Pr(\tau_n\in C_n)\geq1-\alpha\) for every \(n\).

\emph{Oracle comparison.}
The certificate itself and \(\bar q\in\mathcal Q_n(e)\) imply
\[
 \sup_{r'\in\mathcal R}
 \Pr(|Z_{r',\bar q}|>a_n+h_n)
 \leq\sup_{r'\in\mathcal R}
 \mathbb EH_{a_n,h_n}(Z_{r',\bar q})
 \leq\alpha.
\]
Thus \(b_\alpha(\bar q)\leq a_n+h_n\leq b_n+h_n\).

Fix \(\eta\in(0,\alpha)\) and \(s>0\), and put
\(b_\star=b_{\alpha-\eta}(\bar q)\).
Monotone continuity of the tail from above shows that the infimum in the oracle definition is attained, so
\[
 \sup_{r'\in\mathcal R}\Pr(|Z_{r',\bar q}|>b_\star)
 \leq\alpha-\eta.
\]
Since \(H_{b_\star,h_n}\leq{\bf1}\{|z|>b_\star\}\), the same bound holds for its robust ramp risk.  On \(G_n\), every \(q'\in\mathcal Q_n(e)\) satisfies
\(\|q'-\bar q\|_\infty\leq d_{Q,n}\).  The profile-continuity lemma, uniformly in \(r'\), gives
\[
 \sup_{\substack{q'\in\mathcal Q_n(e)\\r'\in\mathcal R}}
 \mathbb EH_{b_\star,h_n}(Z_{r',q'})
 \leq\alpha-\eta+
 \left(\frac{4KB}{c^2}+\frac{2B}{c}\right)\frac{d_{Q,n}}{h_n}.
\]
Here \(d_{Q,n}/h_n\leq\sqrt{d_{Q,n}}\to0\);
also \(D_{\mathrm{tr},n}\to0\) and \(\zeta_n\to0\).
Hence, uniformly over the fixed profile classes, the outward certificate eventually accepts the grid point
\[
 \bar b_\star=2^{-n}\lceil2^n b_\star\rceil
 \leq b_\star+2^{-n}.
\]
Minimality on the corrected grid gives \(a_n\leq\bar b_\star\).  Since
\(h_n\leq x_n+2^{-n}\) and \(x_n\to0\), eventually
\(5x_n+2^{-n}\leq s\), and therefore
\[
 b_n+h_n=a_n+5h_n
 \leq b_\star+s+5\cdot2^{-n}.
\]
This proves the two-sided oracle statement and isolates exactly
\(5(h_n-x_n)\leq5\cdot2^{-n}\) as the dyadic bandwidth charge.

Finally, the cited Li--Zhao result uses a Lindeberg negligible-increment premise for its Gaussian conclusion.  The present class includes, for example,
\(\bar q\equiv c\), \(e_{nt}=c/t\), and a nonzero constant treated residual.
For \(t/n\) bounded away from zero, a treated rare-arm increment has magnitude asymptotic to a nonzero multiple of \(t/n\), while its probability is \(c/t\).
The sum of its truncated second moments over any fixed late interval is bounded away from zero.  Thus the negligible-increment condition fails, although the finite-transfer and atom-safe coverage proof above remains valid.

%%% FIELD proof-ramp-certificate
It suffices to treat the positive sign, since the negative-sign error is its reflection.  Let \(u_t=t/n\).  For \(t>n/2\),
\[
 r^+(u_t)=\frac{u_t-1/2}{2},\qquad
 \frac1n\frac{r^+(u_t)}{e_{nt}}=u_t(u_t-1/2)=:g(u_t).
\]
The finite-population target is
\[
 d_n=\tau_n^+
 =\frac1{2n}\sum_{t=n/2+1}^n(u_t-1/2)
 =\frac1{16}+\frac1{8n}.
\]
Consequently
\[
 \widehat\tau_n-\tau_n
 =\sum_{t=n/2+1}^n A_{nt}g(u_t)-d_n.
\]

For \(i=16,\ldots,31\), place the times with
\(i/32<u_t\leq(i+1)/32\) in cell \(i\).  Couple
\(A_{nt}\sim\operatorname{Bernoulli}(1/(2t))\) as
\({\bf1}\{P_t\geq1\}\), where
\[
 P_t\sim\operatorname{Poisson}(\lambda_t),\qquad
 \lambda_t=-\log(1-1/(2t)).
\]
This has exactly the required Bernoulli law.  Since
\(-\log(1-x)\leq x/(1-x)\),
\[
 \sum_{t\ {\rm in\ cell}\ i}\lambda_t
 \leq\sum_{t\ {\rm in\ cell}\ i}\frac1{2t-1}
 <\frac1{2i}=:\ell_i.
\]
Also \(g\) is increasing on \([1/2,1]\), and its cell maximum is
\[
 G_i=\frac{(i+1)(i-15)}{1024}.
\]
Adding independent Poisson excess counts therefore couples the raw finite sum below
\[
 X=\sum_{i=16}^{31}G_iN_i,\qquad
 N_i\ \hbox{independent},\quad N_i\sim\operatorname{Poisson}(\ell_i).
\]
Because \(d_n<3/8\), the map
\(x\mapsto H_{3/8,\,1/32}(x-d_n)\) is nondecreasing for \(x\geq0\);
and since \(d_n\geq1/16\), it is bounded above by
\(H_{3/8,\,1/32}(x-1/16)\).

For completeness, the directed-integer certificate is the following finite arithmetic identity.  Set \(Q=10^{24}\),
\(\ell_i^+=\lceil Q/(2i)\rceil/Q\), and define the finitely supported rational coefficients
\[
 f_0(0)=1,\qquad
 f_k(x)=\frac1k\sum_{i=16}^{31}\ell_i^+f_{k-1}(x-G_i),
 \quad 1\leq k\leq8.
\]
They are the unnormalised compound-Poisson convolution coefficients through total count eight.  Directed rational evaluation of
\[
 e^{-\sum_i\ell_i^+}
 \sum_{k=0}^8\sum_x f_k(x)
 H_{3/8,\,1/32}(x-1/16)
\]
using the alternating exponential enclosure through degree twelve, and adding
\[
 \frac{0.36^9}{9!\{1-0.36/10\}}
\]
for all counts at least nine, gives the exact outward rational bound
\[
 \frac{37933071186583474830970}{10^{24}}
 =0.037933071186583474830970
 <0.0379330711865835<0.03794.
\]
All multipliers \(G_i\), \(\ell_i^+\), and ramp breakpoints are rational; the recurrence consists only of finitely many directed integer additions, multiplications, and divisions.  Thus this displayed rational is a checkable certificate rather than a simulation result.  Since
\({\bf1}\{|z|>13/32\}\leq H_{3/8,\,1/32}(z)\), the miscoverage probability of the error-radius \(13/32\) interval is below \(0.03794\) for the positive sign and, by reflection, for the negative sign.

%%% FIELD proof-high-hazard-membership
Here \(M_m=0\), so \(B=M\).  The affine profile satisfies
\[
 \bar q_H(0)=1/2=c,\qquad
 1/2\leq\bar q_H(u)\leq100.5\leq K,\qquad
 |\bar q_H(u)-\bar q_H(v)|=100|u-v|\leq L_q|u-v|.
\]
Thus \(\bar q_H\in\mathcal Q\).  For every \(r\in\mathcal R\),
\(Y_{nt}(1)=r(t/n)\), \(Y_{nt}(0)=0\), and \(m_{nt}(a)=0\) are fixed and bounded by \(M\); the regression is predictable, and the treated residual-profile error is identically zero.

Let \(t_0=\lceil n\varepsilon_n\rceil\).  Uniformly for \(t\geq t_0\),
\[
 \frac{1}{2t}+\frac{100}{n}
 \leq\frac{1}{2n\varepsilon_n}+\frac{100}{n}\longrightarrow0.
\]
Hence the cap at \(3/4\) is absent after the cutoff for all sufficiently large \(n\), and then
\[
 t e_{nt}=1/2+100t/n=\bar q_H(t/n)
\]
there.  Thus both stated profile errors vanish.

For every \(t\), the definition of \(e_{nt}\) gives
\(e_{nt}\geq1/(2t)\).  The predictable quadratic variation of the early rare-arm martingale is consequently bounded deterministically by
\[
 \frac1{n^2}\sum_{t<t_0}
 \frac{(1-e_{nt})r(t/n)^2}{e_{nt}}
 \leq\frac{2M^2}{n^2}\sum_{t<t_0}t
 \leq M^2(t_0/n)^2
 =O(M^2\varepsilon_n^2+M^2/n).
\]
Choosing \(v_{0,n}\) equal to this explicit bound and \(\kappa_n=0\) makes the good-event contract hold surely.  The schedule also obeys
\(0<e_{nt}\leq3/4=\rho<1\), so every declared design, overlap, boundedness, predictability, and profile premise holds for all sufficiently large \(n\).

%%% FIELD proof-high-hazard-radius
For a centered Poisson integral, the elementary simple-function calculation used in the finite-transfer proof gives
\[
 \operatorname{Var}(Z_{r,\bar q_H})
 =\int_0^1\frac{u r(u)^2}{1/2+100u}\,du.
\]
On this subfamily \(B=M\), hence \(|r(u)|\leq M\), and
\[
 \operatorname{Var}(Z_{r,\bar q_H})
 \leq M^2\int_0^1\frac{u}{1/2+100u}\,du
 =M^2\frac{100-\frac12\log(201)}{10000}.
\]
Directed Taylor bounds give \(e^5<148.5\) and \(e^{3/10}<1.35\), so
\(e^{5.3}<200.475<201\) and therefore \(\log(201)>5.3\).  It follows that
\[
 \frac{100-\frac12\log(201)}{10000}
 <0.009735.
\]
For any \(b>0\), applying Markov's inequality to \(Z^2\) yields
\[
 \Pr(|Z_{r,\bar q_H}|>b)
 \leq\frac{\operatorname{Var}(Z_{r,\bar q_H})}{b^2}.
\]
Since the exact decimal identity
\[
 0.05(0.442)^2=0.0097682>0.009735
\]
holds, uniformly over \(r\in\mathcal R\),
\[
 \Pr(|Z_{r,\bar q_H}|>0.442M)<0.05.
\]
The definition of the robust oracle radius therefore gives
\(b_{0.05}(\bar q_H)\leq0.442M\).

The inequality is strict.  Choose
\[
 0<\eta<
 0.05-\frac{0.009735}{(0.442)^2}.
\]
Then the same calculation gives
\(b_{0.05-\eta}(\bar q_H)\leq0.442M\).
Apply the oracle upper bound in
\(\mathrm{thm:certified\mbox{-}coverage\mbox{-}oracle}\) with
\(\alpha=0.05\).  For every \(s>0\), on \(G_n\) and for all sufficiently large \(n\),
\[
 b_n+h_n
 \leq b_{0.05-\eta}(\bar q_H)+s+5\cdot2^{-n}
 \leq0.442M+s+5\cdot2^{-n},
\]
as claimed.  Membership of this subfamily in the same primitive class is supplied by
\(\mathrm{prop:high\mbox{-}hazard\mbox{-}membership}\).

%%% FIELD proof-dickman-reduction
Let \(r_0=r(0)\).  Under \(\bar q(u)\equiv c\) and
\(r(u)\equiv r_0\), define
\[
 J_c=\int_{(0,1]}u\,N_c(du),\qquad
 N_c(du)\ \hbox{of intensity }c\,du/u.
\]
The cited generalized-Dickman characterization identifies \(J_c\) with the generalized Dickman law of parameter \(c\).  Direct substitution in the definition of \(Z\) gives
\[
 Z_{r,\bar q}
 =\frac{r_0}{c}\left\{\int_{(0,1]}u\,N_c(du)
       -\int_0^1u\,c\,\frac{du}{u}\right\}
 =\frac{r_0}{c}(J_c-c).
\]
This is precisely the declared centering and scaling of the raw generalized-Dickman jump sum.  If \(r_0=0\), the homogeneous profile is identically zero, so the last display gives \(Z_{r,\bar q}=0\).  Substitution into
\(\mathrm{thm:finite\mbox{-}transfer}\), together with
\(S_n=\widehat\tau_n-\tau_n\), yields
\[
 \Pr\!\left(G_n^c\cup
 \{|\widehat\tau_n-\tau_n|>4h_n\}\right)
 \leq D_{\mathrm{tr},n}.
\]

%%% FIELD stmt-li-zhao-lindeberg
Li and Zhao's Gaussian coverage theorem for the adaptive design-based AIPW estimator uses a Lindeberg-type negligible-increment condition; its conclusion is not asserted for critical rare-arm arrays whose normalized late increments remain nonnegligible.

%%% FIELD stmt-generalized-dickman
For \(\vartheta>0\), the positive infinitely divisible random variable
\[
 W_\vartheta=\int_{(0,1]}u\,N_\vartheta(du),
 \qquad N_\vartheta(du)\text{ having intensity }\vartheta\,du/u,
\]
has the generalized Dickman distribution with parameter \(\vartheta\).

%%% FIELD construction-certificate-corrected
If \(T_n\leq K\), \(\delta_{q,n}>c/4\), the enclosure is empty, or
\(\Delta_n\geq\alpha\), return the full-range flag.  Otherwise let
\(\Gamma_n=\{k2^{-n}:k\in\mathbb N_0\}\).
For each \(b\in\Gamma_n\), \(\mathsf{Cert}_n\) constructs finite rational
Lipschitz outer boxes for \(\mathcal Q_n(e)\) and \(\mathcal R\), truncates the
Poisson integral near zero with a rational cutoff and a second-moment bound,
quantizes time and signed jumps using the common-Poisson Wasserstein bound,
truncates the total count with a certified Poisson-tail bound, performs direct
directed-rational signed-mark convolution, and adds net, quadrature, count, and
rounding allowances totaling at most \(\zeta_n\).  It evaluates the grid points
in increasing order and returns the first \(a_n\in\Gamma_n\) whose outward
upper certificate proves
\[
\sup_{q'\in\mathcal Q_n(e),\,r'\in\mathcal R}
\mathbb E H_{a_n,h_n}(Z_{r',q'})\leq\alpha-\Delta_n.
\]
The grid search may use the uniform second-moment bound
\(\sup_{q',r'}\operatorname{Var}(Z_{r',q'})\leq B^2/c\) as a terminating
large-threshold fallback.  Every comparison and elementary-function
evaluation is made with directed rational interval arithmetic.

%%% FIELD stmt-main-replace
The algorithm \(\mathsf{Cert}_n\) terminates for every certified input. Uniformly over all experiments satisfying the declared assumptions, the reported interval obeys \(\Pr(\tau_n\in C_n)\ge1-\alpha\) for every \(n\), where infeasible budgets trigger \([-2M,2M]\). Moreover, for every \(\eta\in(0,\alpha)\) and \(s>0\), uniformly over the fixed compact profile classes, on \(G_n\) and for all sufficiently large \(n\), \(b_{\alpha}(\bar q)\le b_n+h_n\le b_{\alpha-\eta}(\bar q)+s+5\cdot2^{-n}\). The last term is the complete radius charge for replacing \(x_n\) by its directed-dyadic upper bracket \(h_n\). This covers the declared \(e_{nt}\asymp\bar q(t/n)/t\) profiles that fall outside the Lindeberg condition used for Gaussian coverage in LiZhao2026, Theorem 4(ii).

%%% FIELD stmt-dickman-replace
In the homogeneous-residual special case \(\bar q(u)\equiv c\) and \(r(u)\equiv r(0)\), the raw jump sum underlying \(Z_{r,\bar q}\), after the declared centering and scaling, is the generalized Dickman law with parameter \(c\) studied by BhattacharjeeGoldstein2019. If \(r(0)=0\), then \(Z_{r,\bar q}=0\) and the finite-transfer inequality specializes to \(\Pr(G_n^c\cup\{|\widehat\tau_n-\tau_n|>4h_n\})\le D_{\mathrm{tr},n}\).

%%% FIELD partial-polynomial-certificate
The exhaustive dyadic-grid certificate terminates, is outward-valid, preserves the \(13/32\) late-ramp certificate and the \(0.442M\) high-hazard oracle bound, and yields the proved finite-coverage and oracle-slack theorem.  Moreover, if a fast routine returns any \(\widetilde a_n\in[a_n,a_n+\xi]\) with the stated outward certificate, monotonicity of \(H_{b,h_n}\) preserves coverage and enlarges the reported radius by at most \(\xi\).  What is not proved is a polynomial worst-case bit bound for finding such a threshold.
